% SGA 3, Expose V, Loop 1 English reconstruction.
% Scope: completion of Lemma 7.3 and the page-27 opening of Lemma 7.4.
% Authority: Polo--Gille Exp5-13oct24.pdf, component-local p. 27
% (combined-reader p. 289), re-edited header 277.
% Jacob Reinhold's CC BY 4.0 Markdown is comparison/drafting material only.

因此态射 \(\widetilde d_1\) 在 \(F_1\) 中位于 \(z\) 上的各点处
拟有限且平坦. 由于 \(\widetilde d_1\) 有限型,
由 SGA~1, IV, 6.10 和 EGA~III, 4.4.10 可知,
\(F_1\) 中使 \(\widetilde d_1\) 不同时平坦且拟有限的点所成的集合
\(\Phi\) 在 \(F_1\) 中闭,
并因 \(F_1\) 在 \(X_1\) 中闭而在 \(X_1\) 中闭.%
\footnote{\label{note:sga3-V-40}编者注: 若 \(f:X\to Y\) 局部有限型,
则在其纤维 \(f^{-1}(f(x))\) 中孤立的点 \(x\in X\)
所成的集合在 \(X\) 中开. EGA~III, 4.4.10 在 \(Y\)
局部 Noether 时由 Zariski 主定理推出此点.
对任意 \(Y\), 此点由 Chevalley 半连续性定理
(EGA~IV\(_3\), 13.1.3 和 13.1.4) 得到.
因此, \(f\) 在 \(x\) 处拟有限, 当且仅当它在 \(x\) 处有限型,
且 \(x\) 在 \(f^{-1}(f(x))\) 中孤立.
下面将用到此点; 见\hyperref[note:sga3-V-42]{编者注~(42)}.}
由于 \(d_1\) 固有, \(\widetilde d_1(\Phi)\) 闭,
且由前述内容知它不包含 \(z\). 置
\[
  W=\widetilde d_1(F_1)-\widetilde d_1(\Phi).
\]
\(\widetilde d_1\) 在 \(\widetilde d_1^{-1}(W)\) 上的限制
有限呈示 (由 Noether 假设)、平坦、固有且拟有限.%
\footnote{\label{note:sga3-V-41}编者注: 下面扩充了原文.}
因此, 由 EGA~III, 4.4.2 及 EGA~IV\(_2\), 2.1.12 和 2.4.6,
它有限、局部自由且开. 从而 \(\widetilde d_1(F_1)\) 是 \(z\) 的邻域,
而 \(W\) 是包含在 \(\widetilde d_1(F_1)\) 中,
且使 \(\widetilde d_1\) 在其上同时平坦且拟有限的 \(X_0\) 的最大开集.

我们将在引理~\hyperref[lem:sga3-V-7.4]{7.4} 中看到,
\(\Phi\) 在 \(\widetilde d'_1\) 和 \(\widetilde d'_0\) 下的逆像,
都与 \(F_2\) 中使 \(\widetilde d'_2\) 不同时平坦且拟有限的点所成的集合等同.
于是
\[
\begin{aligned}
  d_0^{-1}(W)
  &=
  \widetilde d'_2(F_2)
  -
  \widetilde d'_2\bigl((\widetilde d'_0)^{-1}(\Phi)\bigr),\\
  d_1^{-1}(W)
  &=
  \widetilde d'_2(F_2)
  -
  \widetilde d'_2\bigl((\widetilde d'_1)^{-1}(\Phi)\bigr)
\end{aligned}
\]
相同. 因而 \(W\) 饱和.

置
\[
  W_1=\widetilde d_1^{-1}(W).
\]
等式 \(d_0^{-1}(W)=d_1^{-1}(W)\) 蕴含
\[
  (\widetilde d'_2)^{-1}\bigl(d_0^{-1}(W)\bigr)
  =
  (\widetilde d'_2)^{-1}\bigl(d_1^{-1}(W)\bigr),
\]
即
\[
  (\widetilde d'_0)^{-1}(W_1)
  =
  (\widetilde d'_1)^{-1}(W_1).
\]
态射 \(\widetilde d_0\) 忠实平坦且拟紧,
因为 \(d_0\) 与 \(d_1\) 一样, 满射、固有且平坦.
此外, 方块
\SGAThreeDiagram[0.30\textwidth]
  {figures/exp5/p277_lemma73_descent_square.png}
  {用于下降 \(W_1\) 的笛卡儿方块.}
  {fig:sga3-V-7.3-descent-square}
  {第 V 篇, 分篇第~27 页 / 合订本第~289 页 / 重版本页眉页码~277}
是笛卡儿图. 可知存在 \(F\) 的开集 \(U\)
(SGA~1, VIII, \SGARefLink{sga3:viii:corollary:4:4}{4.4}), 使得
\[
  W_1=\widetilde d_0^{-1}(U)
\]
此开集 \(U\) 是由 \(X_\ast\) 诱导的以 \(W\) 为底的群胚的拟截面.
因此可取 \(U_z=W\).

还须陈述\hyperref[lem:sga3-V-7.4]{下面的引理}, 其证明是经典的.

\medskip
\phantomsection\label{lem:sga3-V-7.4}%
\noindent\textbf{引理 7.4.}---%
\emph{考虑概形的一个笛卡儿方块}
\SGAThreeDiagram[0.30\textwidth]
  {figures/exp5/p277_lemma74_cartesian_square.png}
  {\hyperref[lem:sga3-V-7.4]{引理~7.4} 中的笛卡儿方块.}
  {fig:sga3-V-7.4-cartesian-square}
  {第 V 篇, 分篇第~27 页 / 合订本第~289 页 / 重版本页眉页码~277}
% The statement continues on component-local p. 28.
% SGA3-V-V2-EOL-PRESERVATION

% Translation coverage: entire assigned component; adjoining components retain the original statement/proof seams.
% Source lines: 103; translated body lines before this audit comment: 99.
% Math preserved: 52 inline formulas, 6 display blocks, 0 equation environments; mathematical prose translated where present.
% Labels: 3/3; hyperrefs: 4/4; SGARefLinks: 1/1; inline targets: 0/0; footnotes: 2/2.
% Figure arguments preserved; all rendered captions and natural-language source locators translated.
% Unresolved TODOs: 0; compilation and visual inspection remain with the volume integrator.
